\section[Rédaction Bourbaki \redac{438} : une identité triangulaire suffit]{Rédaction Bourbaki \redac{438} : une identité triangulaire et un automorphisme font une adjonction\protect\verifie{Redaction438}}\label{sec:bourbaki-438}

La rédaction \redac{438}, \q{Catégories fibrées}, des archives Bourbaki (Archives Henri Poincaré, \ahp{987}) est datée de mars 1965 par sa notice, qui la donne à Claude Chevalley ; elle compte 48 pages. Elle fait partie des textes du livre sur les catégories que Bourbaki n'a pas publié. La démonstration de sa proposition~1, sur les catégories fibrées, a besoin d'une adjonction, et la rédaction isole l'étape dans un lemme~1 (p.~35--36). Sa transcription reste locale pour l'instant ; on n'en cite ici aucun passage.

\emph{Unité, coünité et identités triangulaires.} Une adjonction entre deux foncteurs $F : \mathcal C \to \mathcal D$ et $G : \mathcal D \to \mathcal C$ peut se donner par deux transformations naturelles, une unité $\eta : \mathrm{Id} \to GF$ et une coünité $\varepsilon : FG \to \mathrm{Id}$, soumises à deux identités triangulaires :
\[
F X \xrightarrow{F\eta_X} FGF X \xrightarrow{\varepsilon_{FX}} F X \quad\text{et}\quad G Y \xrightarrow{\eta_{GY}} GFG Y \xrightarrow{G\varepsilon_Y} G Y
\]
sont les identités. L'exemple type est celui du module libre : $F$ envoie un ensemble sur le module libre qu'il engendre, $G$ oublie la structure de module, $\eta$ envoie un élément sur le vecteur de base correspondant, et $\varepsilon$ évalue une combinaison formelle d'éléments d'un module. Les deux identités disent que ces deux opérations se compensent, d'un côté comme de l'autre.

\emph{Le lemme.} Dans la rédaction, les rôles de $F$ et $G$ sont tenus par deux foncteurs notés $R$ et $S$. On suppose qu'une des deux identités triangulaires a lieu, et, de l'autre composé $G\varepsilon_Y \circ \eta_{GY}$, seulement qu'il est un automorphisme de $GY$ pour tout $Y$ (figure~\ref{fig:bourbaki-438}). Le lemme conclut que ce composé est l'identité, de sorte que $\eta$ et $\varepsilon$ sont l'unité et la coünité d'une adjonction.

\begin{figure}[htbp]
\centering
\begin{tikzcd}
FX \arrow[r, "F\eta_X"] \arrow[dr, "\mathrm{id}_{FX}"'] & FGFX \arrow[d, "\varepsilon_{FX}"] \\
& FX
\end{tikzcd}
\qquad\qquad
\begin{tikzcd}
GY \arrow[r, "\eta_{GY}"] \arrow[dr, dashed, "e_Y"'] & GFGY \arrow[d, "G\varepsilon_Y"] \\
& GY
\end{tikzcd}
\caption{Les deux identités triangulaires. Le lemme suppose que le triangle de gauche commute pour tout $X$, et seulement que le composé $e_Y = G\varepsilon_Y \circ \eta_{GY}$ du triangle de droite (en tirets) est un isomorphisme ; il en tire $e_Y \circ e_Y = e_Y$, donc $e_Y = \mathrm{id}_{GY}$, et le triangle de droite commute aussi.}
\label{fig:bourbaki-438}
\end{figure}

\begin{encadre}
\begin{enonce}[Bourbaki, rédaction \redac{438}, lemme~1, p.~35]
Soient $F : \mathcal C \to \mathcal D$, $G : \mathcal D \to \mathcal C$, $\eta : \mathrm{Id} \to GF$ et $\varepsilon : FG \to \mathrm{Id}$. Si $\varepsilon_{FX} \circ F\eta_X = \mathrm{id}_{FX}$ pour tout $X$, et si $e_Y = G\varepsilon_Y \circ \eta_{GY}$ est un isomorphisme pour tout $Y$, alors $e_Y = \mathrm{id}_{GY}$ pour tout $Y$, et $(\eta, \varepsilon)$ définit une adjonction $F \dashv G$. Il en va de même en échangeant les rôles des deux identités triangulaires.
\end{enonce}
\end{encadre}

\emph{L'idée.} Un projecteur inversible est l'identité. Il suffit donc de voir que $e_Y$ est idempotent, $e_Y \circ e_Y = e_Y$ : on compose alors par l'inverse de $e_Y$, et il reste $e_Y = \mathrm{id}$. La rédaction procède autrement, en comparant les deux applications entre ensembles de morphismes qu'induisent $\eta$ et $\varepsilon$. La preuve vérifiée est plus courte.

\begin{proof}
On fixe $Y$ et on écrit $e = G\varepsilon_Y \circ \eta_{GY}$. Dans $e \circ e = G\varepsilon_Y \circ \eta_{GY} \circ G\varepsilon_Y \circ \eta_{GY}$, on fait passer les deux flèches du milieu l'une par-dessus l'autre, par naturalité de $\eta$ relativement à $G\varepsilon_Y$ :
\[
\eta_{GY} \circ G\varepsilon_Y = GFG\varepsilon_Y \circ \eta_{GFGY}.
\]
On obtient $G\varepsilon_Y \circ GFG\varepsilon_Y \circ \eta_{GFGY} \circ \eta_{GY}$. La naturalité de $\varepsilon$ relativement à $\varepsilon_Y$ donne $\varepsilon_Y \circ FG\varepsilon_Y = \varepsilon_Y \circ \varepsilon_{FGY}$, et celle de $\eta$ relativement à $\eta_{GY}$ donne $\eta_{GFGY} \circ \eta_{GY} = GF\eta_{GY} \circ \eta_{GY}$. Il vient
\[
e \circ e = G\varepsilon_Y \circ G\bigl(\varepsilon_{FGY} \circ F\eta_{GY}\bigr) \circ \eta_{GY}.
\]
La parenthèse est l'identité triangulaire supposée, au point $X = GY$ ; elle vaut l'identité, et il reste $e \circ e = e$. Comme $e$ est inversible, $e = \mathrm{id}$, et les deux identités triangulaires ont lieu. Le cas dual se traite de la même façon, avec le composé $\varepsilon_{FX} \circ F\eta_X$.
\end{proof}

\noindent\emph{Vérification en Lean.} L'énoncé encadré est \texttt{rightTriangle\_of\_leftTriangle}, appuyé sur \texttt{rightComposite\_idem} (le composé est idempotent) et \texttt{eq\_id\_of\_isIso\_of\_idem} (un idempotent inversible est l'identité). Le dual est \texttt{leftTriangle\_of\_rightTriangle}. Les adjonctions elles-mêmes sont \texttt{adjunctionOfLeftTriangle} et \texttt{adjunctionOfRightTriangle}, construites par \texttt{Adjunction.\allowbreak{}mkOfUnitCounit} de mathlib, dans \href{\leanblobmain/Redaction438.lean}{\texttt{Redaction438.\allowbreak{}lean}}.

Ne sont pas démontrées ici : la proposition~1 de la rédaction, sur les catégories fibrées et le produit fibré canonique, dont le lemme est une étape.

\litterature{l'énoncé est probablement le lemme~A1.1.1 de l'\emph{Elephant} de Johnstone \cite{johnstone2002} ; la comparaison avec le livre n'a pas été faite.}
